% Chapter 3: velocity sets

\chapter{Velocity sets}

\label{Chapter3}

The module \code{fluxlb.core.lattices} provides the discrete velocity sets. Five lattices
are implemented: D2Q5, D2Q9, D2Q37, D3Q19 and D3Q27.

%----------------------------------------------------------------------------------------

\section{Theory}

\begin{maintainer}[discrete velocity sets and the quadrature view]
Derive or motivate the discrete velocity sets: the Hermite expansion of the
Maxwell--Boltzmann distribution, Gauss--Hermite quadrature on a lattice, and why the weights
$w_i$ and sound speed $\cs^2$ follow from requiring the discrete moments to match the
continuous ones up to a chosen order. State which moment orders are needed to recover
(i) advection--diffusion, (ii) isothermal Navier--Stokes and (iii) the energy equation, and
connect each to the \code{isotropy\_order} attribute of \cref{sec:lattice-impl}. Cite the
sources you have read \parencite{Philippi2006,Sbragaglia2009,Kruger2017}.
\end{maintainer}

%----------------------------------------------------------------------------------------

\section{Implementation}\label{sec:lattice-impl}

A lattice is an \code{nn.Module} that holds constants and no computation. The base class
\code{Lattice} registers four non-persistent buffers and one plain attribute, listed in
\cref{tab:lattice-contents}. Class attributes \code{Q}, \code{D} and \code{isotropy\_order}
are integers.

\begin{table}[h]
\centering\small
\begin{tabular}{@{}llll@{}}
\toprule
Name & Shape & Dtype & Role \\
\midrule
\code{c}   & $[Q, D]$ & compute dtype & velocities, used in the momentum sum \\
\code{w}   & $[Q]$    & compute dtype & weights \\
\code{opp} & $[Q]$    & \code{int64}  & $\bar{i}$, for bounce-back; indexes $f$ on the same device \\
\code{cs2} & scalar   & compute dtype & $\cs^2$ \\
\code{shifts} & $Q$ tuples & Python \code{int} & velocities for \code{torch.roll}; never on the device \\
\bottomrule
\end{tabular}
\caption{Contents of a \code{Lattice}.}
\label{tab:lattice-contents}
\end{table}

\paragraph{Why buffers.}
\code{register\_buffer} makes the constants follow \code{.to(device, dtype)} together with
the solver, keeps them out of \code{parameters()} so no optimiser ever sees them, and casts
only floating-point buffers on a dtype change, so \code{opp} stays \code{int64}. They are
non-persistent: derived data stays out of checkpoints and cannot be loaded back with a stale
dtype.

\paragraph{Why \code{shifts} is not a buffer.}
\code{torch.roll} takes Python integers. Reading them from a GPU tensor inside the time loop
would force a device-to-host synchronisation every step.

\paragraph{Dtype.}
The constructor takes \code{dtype} (default \code{torch.float32}). Weights are written as
exact Python fractions and cast once, so an fp64 lattice has correctly rounded constants
rather than upcast fp32 ones. Building in fp64 is how the CPU regression tests will run.

\paragraph{Validation at construction.}
The base raises \code{ValueError} if the shapes of \code{c}, \code{w} or \code{opp} disagree
with \code{Q} and \code{D}, or if \code{c} is not integer-valued. The second check exists
because \code{shifts} is built with \code{int()}; a scaled stencil would otherwise be
truncated silently and stream the wrong way.

\paragraph{Declared isotropy order.}
Each lattice declares the highest even order at which its weighted velocity moments are
isotropic. The tests check exactly the orders a lattice claims, and one negative test
confirms that D2Q5 genuinely fails at fourth order so the attribute cannot quietly be
overstated.

%----------------------------------------------------------------------------------------

\section{Lattice catalogue}

All lattices use integer velocities and stream with \code{torch.roll}. Within every lattice
the rest velocity is at index~0. \Cref{tab:catalogue} summarises them; ordering conventions
follow.

\begin{table}[h]
\centering\small
\begin{tabular}{@{}lcccccl@{}}
\toprule
Lattice & $Q$ & $D$ & Isotropy order & $\cs^2$ & Distinct weights & Intended use \\
\midrule
D2Q5  & 5  & 2 & 2 & $1/3$   & 2 & scalar transport only \\
D2Q9  & 9  & 2 & 4 & $1/3$   & 3 & isothermal flow \\
D2Q37 & 37 & 2 & 8 & $1/r^2$ & 8 & thermal / compressible \\
D3Q19 & 19 & 3 & 4 & $1/3$   & 3 & isothermal flow \\
D3Q27 & 27 & 3 & 4 & $1/3$   & 4 & isothermal flow, full stencil \\
\bottomrule
\end{tabular}
\caption{Implemented lattices. $r = 1.19697977039307435897239$ for D2Q37.}
\label{tab:catalogue}
\end{table}

\subsection{D2Q5}

Rest, then the four axis directions anticlockwise from $+x$:
\[
(0,0),\ (1,0),(0,1),(-1,0),(0,-1).
\]
Weights $1/3$ (rest) and $1/6$ (axis); $\code{opp} = [0,3,4,1,2]$.

\begin{verified}
Normalised; second moment equals $\tfrac{1}{3}\delta_{\alpha\beta}$; fourth-order isotropy
fails with maximum error $0.11$. Declared \code{isotropy\_order = 2}.
\end{verified}

\begin{maintainer}[why D2Q5 is restricted to scalar transport]
State what the missing fourth-order isotropy means for the viscous stress and hence why the
lattice can carry a temperature or concentration field coupled to a D2Q9 flow but not the
flow itself.
\end{maintainer}

\subsection{D2Q9}

Rest, then the four axis directions anticlockwise from $+x$, then the four diagonals
anticlockwise from $(1,1)$:
\[
(0,0),\ (1,0),(0,1),(-1,0),(0,-1),\ (1,1),(-1,1),(-1,-1),(1,-1).
\]
Weights $4/9$, $1/9$ (axis), $1/36$ (diagonal). $\code{opp} = [0,3,4,1,2,7,8,5,6]$.

\begin{verified}
Normalised to rounding; first and third moments vanish; second moment
$\tfrac13\delta_{\alpha\beta}$; fourth-order isotropy holds to $3\times10^{-8}$ in fp32 and
$6\times10^{-17}$ in fp64; sixth order fails. Declared \code{isotropy\_order = 4}.
\end{verified}

\subsection{D2Q37}

The space-filling form of the D2V37 quadrature of \textcite{Philippi2006}: integer
velocities on eight shells with one weight per shell, ordered shell by shell with opposite
pairs adjacent within a shell. \code{opp} is derived from the velocity list rather than
typed. \Cref{tab:d2q37} lists the shells.

\begin{table}[h]
\centering\small
\begin{tabular}{@{}clcl@{}}
\toprule
Shell & Representative $\vc$ & Multiplicity & Weight \\
\midrule
0 & $(0,0)$ & 1 & 0.23315066913235250228650 \\
1 & $(1,0)$ & 4 & 0.10730609154221900241246 \\
2 & $(1,1)$ & 4 & 0.05766785988879488203006 \\
3 & $(2,0)$ & 4 & 0.01420821615845075026469 \\
4 & $(2,1)$ & 8 & 0.00535304900051377523273 \\
5 & $(2,2)$ & 4 & 0.00101193759267357547541 \\
6 & $(3,0)$ & 4 & 0.00024530102775771734547 \\
7 & $(3,1)$ & 8 & 0.00028341425299419821740 \\
\bottomrule
\end{tabular}
\caption{D2Q37 shells. Within a shell, opposite pairs are adjacent; see the class docstring
for the full order. $\cs^2 = 1/r^2 = 0.697953322019683$.}
\label{tab:d2q37}
\end{table}

\begin{verified}
The eight weights were re-derived independently by solving the moment constraints of orders
0 to 8 for the eight shell weights given $r$; the solution agrees with the tabulated values to
$3\times10^{-15}$ and the constraints are satisfied to $7\times10^{-15}$. Odd moments vanish
and even moments are isotropic through eighth order. Declared \code{isotropy\_order = 8}.
\end{verified}

\begin{maintainer}[D2Q37: scaling, reference temperature and thermal use]
Confirm $r$ and the weights against \textcite{Philippi2006} and \textcite{Sbragaglia2009}.
Explain the relation between the quadrature abscissae and the integer lattice, why the
reference temperature $T_0 = 1/r^2$ appears as $\cs^2$, and which moments the thermal
equilibrium needs.
\end{maintainer}

\subsection{D3Q19}

Rest; six axis directions; twelve edge (face-diagonal) directions. Opposite pairs are
adjacent throughout:
\[
+x,-x,+y,-y,+z,-z,\quad\text{then}\quad (1,1,0),(-1,-1,0),(1,-1,0),(-1,1,0)
\]
and likewise in the $xz$ and $yz$ planes. Weights $1/3$, $1/18$ (axis), $1/36$ (edge).

\begin{verified}
Normalised; odd moments vanish; second moment $\tfrac13\delta_{\alpha\beta}$; fourth-order
isotropy holds to fp64 rounding; sixth order fails (error $0.22$). Declared
\code{isotropy\_order = 4}.
\end{verified}

\subsection{D3Q27}

The D3Q19 ordering followed by the eight corner (body-diagonal) directions, again in
opposite pairs:
\begin{gather*}
(1,1,1),(-1,-1,-1),\quad (1,1,-1),(-1,-1,1),\\
(1,-1,1),(-1,1,-1),\quad (-1,1,1),(1,-1,-1).
\end{gather*}
The first nineteen entries coincide with D3Q19. Weights $8/27$, $2/27$ (axis), $1/54$
(edge), $1/216$ (corner).

\begin{verified}
Same identities as D3Q19 hold; sixth order fails (error $0.22$). Declared
\code{isotropy\_order = 4}.
\end{verified}

\begin{maintainer}[D3Q19 versus D3Q27]
State what the corner velocities add (which third-order moments, and the consequence for
Galilean invariance and stability) and when the extra eight populations are worth the memory
on a \SI{16}{GB} card.
\end{maintainer}

%----------------------------------------------------------------------------------------

\section{Using a lattice}

\begin{lstlisting}[style=py]
import torch
from fluxlb.core.gpu import get_device, accumulation_dtype
from fluxlb.core.lattices import D2Q9

device = get_device()
lat = D2Q9().to(device)          # buffers move; shifts stay Python ints

f = torch.rand(lat.Q, 128, 128, device=device)
acc = accumulation_dtype(device)
rho = f.sum(dim=0, dtype=acc).to(f.dtype)
mom = torch.einsum("qa,q...->a...", lat.c.to(acc), f.to(acc)).to(f.dtype)
\end{lstlisting}

Streaming is not yet implemented; when it is, each population $f_i$ rolls by
\code{lat.shifts[i]} along the spatial axes, and bounce-back reads $f_{\bar{i}}$ via
\code{f[lat.opp]}.
